%!TEX root = ./../../thesis-trier.tex

\IN{circular automaton}
An automaton $\semiautomat{A}{Q}{\Sigma}{\delta}$ is \emph{circular}
if there exists a letter $b \in \Sigma$ 
that cyclically permutes the states, i.e., $Q = \{ \delta(q, b^i) \mid i \ge 0 \}$
and $\delta(q, b^{|Q|}) = q$ for some state $q \in Q$ (and hence all states $q \in Q$).
The automata from the \Cerny family (see Figure~\ref{intro:fig:cerny}) are circular.
Dubuc~\cite{Dubuc98} confirmed the \Cerny conjecture for circular automata.

In \cref{syncmax:classes:thm:circular}, we first give a condition implying complete reachability,
complete distinguishability, and bounds on the lengths of shortest words reaching non-empty subsets.
This result was stated by Don~\cite[Prop.~15]{Don16} for the special case of a single
letter of defect one and without concluding complete distinguishability.
We then derive several consequences and give examples.

\begin{theorem} \label{syncmax:classes:thm:circular} \label{sync_max:prop:don}
	Let $\semiautomat{A}{Q}{\Sigma}{\delta}$ be an $n$-state automaton with $\Sigma = \{a_1, \ldots, a_k, b\}$
	where $b \in \Sigma$ cyclically permutes the states and
	$a_1, \ldots, a_k \in \Sigma$ are letters of defect one
	with $0 < d_1, \ldots, d_k < n$	such that
	\begin{enumerate}
	\item for all $i \in \{1,\ldots,k\}$ we have $\dupl{a_i} = \delta(\excl{a_i}, b^{d_i})$ and
	
	\item the greatest common divisor of $d_1, \ldots, d_k$ and $n$ is one.
	\end{enumerate}
	Then $\mathcal A$ is completely distinguishable,
	every non-empty subset of size~$r$ is reachable by a word of length $n(n - r)$,
	and a shortest synchronizing word has length at most $(n - 1)^2$.
\end{theorem}
\begin{proof}
	Consider the digraph $\Gamma_{1,n}(\mathcal A)$ defined on page~\pageref{def:Gamma1i-graph}.
	First, we show that the group of units~$G$ of $\TM{\mathcal A}$ is contained in the 
	automorphism group of $\Gamma_{1,n}(\mathcal A)$.
	Since~$G$ is generated 
	by the action of the letter~$b$, this follows if for every edge $(p,q)$ in $\Gamma_{1,n}(\mathcal A)$
	also the edge $(\delta(p, b), \delta(q, b))$ is in $\Gamma_{1,n}(\mathcal A)$.	
	Let $(p, q)$ be an edge in $\Gamma_{1,n}(\mathcal A)$.
	Then it is induced by a word of the form $a_i b^{j}$
	for some $i \in \{1,\ldots,k\}$ and $j \in \{0,1,\ldots,n-1\}$.
	For the word $u = a_i b^{(j+1) \bmod n}$
	we have $\excl{u} = \{ \delta(p, b)\}$
	and $\dupl{u} = \{ \delta(q, b) \}$.
	Thus, $u$ induces the edge $(\delta(p, b), \delta(q, b))$ in $\Gamma_{1,n}(\mathcal A)$.
		
	Next, assume $\Gamma_{1,n}(\mathcal A)$ is not strongly connected.
	An automorphism maps a path to another path.
	Thus, since the elements from~$G$ act as automorphisms, 
	the weakly connected components form the congruence classes of
	a $G$-congruence on~$Q$.
	Then, by \cref{lem:congruence_on_cyclic_group} there exists a divisor $d$ of~$n$ such that
	two states $p,q \in Q$ are in a common weakly connected component
	if and only if $\delta(p, b^j) = q$
	for some multiple~$0 \le j < n$ of~$d$.
	Note that since~$d$ divides~$n$, the numbers
	$j$ and $n - j$ are either both divisible by~$d$
	or neither is divisible by~$d$.	
	For $i \in \{1,\ldots,k\}$,
	let $\excl{a_i} = \{p_i\}$
	and $\dupl{a_i} = \{q_i\}$.
	Since $(p_i, q_i)$ is an edge in $\Gamma_{1,n}(\mathcal A)$,
	the states~$p_i$ and~$q_i$ are in the same weakly connected component.
	Thus $\delta(p_i, b^j) = q_i$ for some multiple $0 \le j < n$ of~$d$.
	Since $\delta(p_i, b^{d_i}) = q_i$,
	we find $j = d_i$.
	This yields that~$d$ is a common divisor
	of $d_1, \ldots, d_k$ and~$n$.
	By assumption, we find $d = 1$. This implies that all states
	are congruent, i.e., $\Gamma_{1,n}(\mathcal A)$ is weakly connected.
	Since the group $G$ is transitive, $\Gamma_{1,n}(\mathcal A)$ is vertex transitive.
	Thus, by \cref{prelim:thm:vertex_transitive_finite_graph}, it is strongly connected.
	Then, by \cref{syncmax:thm:Gamma1-SCC-subgraph}, 
	$\mathcal A$ is completely reachable,
	each non-empty subset of size $r$ is reachable by a word of length $n(n - r)$,
	and a shortest synchronizing word has length at most $1 + n(n-2) = (n-1)^2$.
	By \cref{thm:defectatmostone}, we find that~$\mathcal A$ is completely distinguishable.
\end{proof}

The following corollary was stated in~\cite[Thm.~13]{HoffmannCircularExtended} in the special case $k = 1$.

\begin{corollary}  \label{sync_max:thm:circular_sync_max} \label{sync_max:cor:circular}
	Let $\semiautomat{A}{Q}{\Sigma}{\delta}$ be an $n$-state automaton where $b \in \Sigma$ cyclically permutes the states.
	Suppose $a_1, \ldots, a_k \in \Sigma$ are letters of defect one
	and $0 < d_1, \ldots, d_k < n$
	are such that for all $i \in \{1,\ldots,k\}$ we have $\delta(q_i, a_i) = \delta(q_i, b^{d_i} a_i)$
	for some $q_i \in Q$.
	If the greatest common divisor of $d_1, \ldots, d_k$ and $n$
	is one, then $\mathcal A$ is completely reachable and completely distinguishable. % (and so sync-maximal).	
\end{corollary}
\begin{proof}
	Since $\TM{\mathcal A}$ contains the transitive permutation group $G = \langle \{ \fword{\delta}{u} \} \rangle_G$,
	by \cref{sync_words:lem:idempotent}, for each $i \in \{1,\ldots,k\}$
	there exist simple idempotent words $u_i$
	mapping $q_i$ to $\delta(q_i, b^{d_i})$.
	We have $\excl{u_i} = \{q_i\}$
	and $\dupl{u_i} = \{ \delta(q_i, b^{d_i}) \}$.
	Considering a new automaton with state set~$Q$ with the letter~$b$
	acting in the same way and where we consider the words $u_1, \ldots, u_k$ as letters,
	from \cref{syncmax:classes:thm:circular} 
	we get that
	this automaton is completely reachable and completely distinguishable.
	This also yields that~$\mathcal A$ is completely reachable and completely distinguishable.
\end{proof}

We will apply \cref{sync_max:thm:circular_sync_max} 
in \cref{sync_max:sec:slowly_sync} (see \cref{sync_max:thm:families}) to show
sync-maximality for a variety of automata families (among them one for which this was 
previously only conjectured).

Since in the proof of \cref{sync_max:thm:circular_sync_max} we have applied \cref{syncmax:classes:thm:circular} 
when considering certain words as letters, the bounds from \cref{sync_max:prop:don}
do not hold (but they transfer when multiplied by the maximum of the lengths
of these words).

Recall that, by \cref{cor:char-bin-compl-reach}, for binary automata with at least three states,
complete reachability is equivalent to sync-maximality.
However, Complete distinguishability in itself does not imply complete reachability
for binary circular automata. For example, \cref{sync_max:fig:compl_dist_not_compl_reach:1} 
gave a binary circular automaton that is completely distinguishable, but not completely reachable.

\Cref{sync_max:cor:circular} can be used to construct sync-maximal circular automata. We give a few
examples in the following statement. See also \cref{syncmax:fig:binary-circular-sync-maximal} which
depicts some cases from this statement.

\begin{figure}[htbp]
	\centering
	\begin{tikzpicture}
		
		\matrix (m) [matrix of nodes,
		column sep=1.5cm,
		row sep=0.01cm] {
			% nodes={anchor=center}
			
			\begin{tikzpicture}[inner sep=2pt, minimum size=0pt, scale=0.65, transform shape, ->, >=stealth]
				\foreach \i in {0,...,6} {
					\node[state] (q\i) at ({360/7 * \i}:3cm) {$\i$};
				}
				
				\foreach \i [evaluate=\i as \j using {int(mod(\i+1,7))}] in {0,...,6} {
					\path[->] (q\i) edge[bend right=15, dashed] (q\j);
				}
				
				\path[->] (q0) edge (q2);
				\draw (q3) edge (q0); % choose odd i=3				
			\end{tikzpicture}
			
			&
			%	\begin{tikzpicture}[scale=0.75, transform shape, ->, >=stealth, shorten >=1pt, shorten <=1pt]
				\begin{tikzpicture}[scale=0.65, transform shape, ->, >=stealth, node distance=2.2cm, on grid, auto]
					\foreach \i in {0,...,6} {
						\node[state] (q\i) at ({360/7 * \i}:3cm) {$\i$};
					}
						
					\draw (q0) edge[dashed] (q1);
					\draw (q1) edge[bend left=8, dashed] (q2);
					\draw (q2) edge[bend left=8, dashed] (q3);
					\draw (q3) edge[dashed] (q4);
					\draw (q4) edge[bend left=8, dashed] (q5);
					\draw (q5) edge[bend left=8, dashed] (q6);
					\draw (q6) edge[bend left=8, dashed] (q0);
					
					\draw (q0) edge[bend left=20] (q3);
					\draw (q3) edge[bend right=10] (q1); % odd i=3
					
					\foreach \i in {1,2,4,5,6} {
						\pgfmathtruncatemacro{\j}{mod(\i+1,7)}
						\draw (q\i) edge[bend right=8]  (q\j);
					}
				\end{tikzpicture}	
			\\			
			{Automaton as in \crefitem{prop:more_families}{prop:more_families:1}}
			&
			{Automaton as in \crefitem{prop:more_families}{prop:more_families:2}}
			\\
			
			\begin{tikzpicture}[inner sep=5pt, minimum size=0pt, scale=0.65, transform shape, ->, >=stealth, node distance=2.2cm, on grid, auto]
				\foreach \i in {0,...,7} {
					\node[state] (q\i) at ({360/8 * \i}:3cm) {$\i$};
				}
				
				\foreach \i [evaluate=\i as \j using {int(mod(\i+1,8))}] in {0,...,7} {
					\draw (q\i) edge[bend right=15, dashed] (q\j);
				}
				
				\draw (q0) edge (q3);
				\draw (q5) edge (q0);
				
			\end{tikzpicture}
			&
			\begin{tikzpicture}[inner sep=5pt, minimum size=0pt, scale=0.65, transform shape, ->, >=stealth, node distance=2.2cm, on grid, auto]
				\foreach \i in {0,...,7} {
					\node[state] (q\i) at ({360/8 * \i}:3cm) {$\i$};
				}
				
				\foreach \i [evaluate=\i as \j using {int(mod(\i+1,8))}] in {0,...,7} {
					\draw (q\i) edge[bend right=15, dashed] (q\j);
				}
				
				\draw (q0) edge[bend left=10]  (q1);	
				
				\draw (q5) edge[bend left=15]  (q6);
				\draw (q6) edge[bend left=15]  (q7);
				\draw (q7) edge[bend right=30] (q5);
			\end{tikzpicture}
			\\
			
			{Automaton as in \crefitem{prop:more_families}{prop:more_families:3}}
			&
			{Automaton as in \crefitem{prop:more_families}{prop:more_families:4}}
			\\
		};		
	\end{tikzpicture}
	\caption[Example sync-maximal circular automata]{Solid and dashed edges show the actions 
		of~$a$ and, resp., $b$; self-loop omitted.}
	\label{syncmax:fig:binary-circular-sync-maximal}
\end{figure}




\begin{proposition}
	\label{prop:more_families}
	Let $\semiautomat{A}{Q}{\{a,b\}}{\delta}$
	with $Q = [n]$ and $\delta(i, b) = (i + 1) \bmod n$.
	Then~$\mathcal A$ is sync-maximal if~$a$ is defined in the following way:
	\begin{enumerate}
		\item \label{prop:more_families:1} $\delta(0, a) = 2$, precisely one odd $i \in \{1,3,\ldots,n-2\}$
		such that $\delta(i, a) = 0$ and $\delta(j, a) = j$ for $j \in [n] \setminus\{0,i\}$ with $n > 2$ being odd,
		
		\item \label{prop:more_families:2} $\delta(0, a) = 3$, precisely one odd $i \in \{1,3,\ldots,n-2\}$
		such that $\delta(i, a) = 1$ and $\delta(j, a) = (j+1) \bmod n$ for $j \in [n] \setminus\{0,i\}$ with $n > 2$ being odd,
		
		\item \label{prop:more_families:3} $\delta(0, a) = 3$, precisely one $i \in \{1,3,\ldots,n-2\}$
		not divisible by $3$ such that $\delta(i, a) = 0$ and $\delta(j, a) = j$ for $j \in [n] \setminus\{0,i\}$ where $n \not\equiv 0 \pmod{3}$ with $n > 3$,
		
		\item \label{prop:more_families:4} $\delta(0,a) = 1$, $\delta(i, a) = i$ for $i \in \{1,\ldots,\lfloor n / 2 \rfloor\}$,
		and $a$ acts injectively on $\{\lfloor n/2 \rfloor + 1, \ldots, n - 1 \}$
		i.e., $\delta(j, a) = \delta(j', a)$ implies $j = j'$ for $j,j' \in \{\lfloor n/2 \rfloor + 1, \ldots, n - 1 \}$.
	\end{enumerate}
\end{proposition}
\begin{proof}
	All these automata are circular with~$a$ having defect one.
	Then the claims follow from \cref{sync_max:cor:circular} due to the following 
	reasons:
	
	\ref{prop:more_families:1}: We have $\delta(0, a) = \delta(0, b^2a)$ and $n$ is odd.
	
	\ref{prop:more_families:2}: We have $\delta(0, a) = \delta(0, b^2a)$ and $n$ is odd.
	
	\ref{prop:more_families:3}: We have $\delta(0, a) = \delta(0, b^3a)$ and $n$ is not a multiple of~$3$.
	
	\ref{prop:more_families:4}: We have $\delta(0, a) = \delta(0, ba)$.
\end{proof}

As the next example shows, \cref{syncmax:classes:thm:circular} applies to a wider 
class of automata than \cref{sync_max:thm:circular_sync_max}.

\begin{example} 
	\label{ex:thm_not_usable}
	Consider the automaton from \cref{ex:I_vs_Gamma1} (see \cref{fig:thm_not_usable}).
	Let~$w$ be a word of defect one. As can be read off
	from  Table~\ref{tab:rank_3_words},
	for $d > 0$ and $q \in [4]$, 
	if 
	$\applytrans{\delta}{q}{w} = \applytrans{\delta}{\applytrans{\delta}{q}{b^d}}{w}$,
	then $d$ is a multiple of two.
	So \cref{sync_max:thm:circular_sync_max}
	does not apply here.
	We have  $\excl{a} = \{0\}$, $\dupl{a} = \{1\}$,
	and $\delta(0, a) = 1 = \delta(0, b)$.
	Thus \cref{sync_max:prop:don} yields that the automaton
	is completely reachable (which also can be seen in \cref{fig:thm_not_usable}).
\end{example}

Next, we give examples of completely reachable binary automata for which complete reachability does not follow directly from \cref{sync_max:prop:don}. In \cref{syncmax:ex:circular-compl-reach-one}, by considering certain words as letters, it becomes possible to apply \cref{sync_max:prop:don}. However, it is unclear whether such an argument can also be applied in \cref{syncmax:ex:circular-compl-reach-two}.

\begin{example} \label{syncmax:ex:circular-compl-reach-one}
	Let $\semiautomat{A}{Q}{\Sigma}{\delta}$
	where $Q = [n]$ for some $n \ge 3$
	and $\delta(q, b) = (q + 1) \bmod n$ for $q \in Q$,
	$\delta(0, a) = 2$, and $\delta(q, a) = q + 1$
	for $q \in \{1,2,\ldots,n-2\}$
	and $\delta(n-1,a) = 1$.
	For even~$n$, \cref{sync_max:prop:don}
	is not applicable.
	However, we have	
	$\delta(0, a^{n}) = 1$,	
	$a^{n}$ has defect one,
	and $0$ and~$1$ are the excluded and duplicate
	state of~$a^{n}$, respectively. Thus, regarding $a^{n}$
	as an alphabet letter, from \cref{sync_max:prop:don}
	we get that~$\mathcal A$ is completely reachable.
\end{example}

\begin{example} \label{syncmax:ex:circular-compl-reach-two}
	Let $\semiautomat{A}{Q}{\Sigma}{\delta}$
	where $Q = [9]$, $\delta(q, b) = (q + 1) \bmod 9$ for $q \in Q$
	and $\delta(0, a) = 3$, $\delta(a, 6) = 7$, $\delta(7, a) = 6$,
	and $\delta(q, a) = q$ for $q \in Q \setminus\{0,6,7\}$, see \cref{syncmax:fig:example-cr-circular}.
	Next, we show that~$\mathcal A$ is completely reachable.
	Since $\excl{a} = \delta(\dupl{a}, b^3)$,
	this does not follow from \cref{sync_max:prop:don},
	and it is also not clear to us how to find words
	acting in a way to apply \cref{sync_max:prop:don}.
	Next, let $S \subseteq Q$ be a non-empty subset of~$Q$.
	We show, by induction on $|Q \setminus S|$,
	that~$S$ is reachable in $\mathcal A$.
	The claim holds for $S = Q$.
	Next, let $S \ne Q$.	
	
	\begin{enumerate}
	\item Suppose we have $\delta(p, b^3) \in S$ for some $p \in Q \setminus S$.
	
	In this case, there exists $i \ge 0$ such that
	$0 \notin \delta^{-1}(S, b^i)$ and $3 \in \delta^{-1}(S, b^i)$.
	This gives $|\delta^{-1}(S, ab^i)| = |S| + 1$
	and $\delta( \delta^{-1}(S, ab^i), ab^i ) = S$.
	Inductively, we have $\delta(Q, u) = \delta^{-1}(S, ab^i)$
	for some $u \in \Sigma^*$.
	Then $\delta(Q, uab^i) = S$.	
	
	\item Otherwise, each residue class modulo three is either contained in $S$ or in its complement.
		
	Let $q \in Q \setminus S$ be such that $\delta(q, b) \in S$.
	Then we have $6 \notin \delta^{-1}(S, b^i)$ and $7 \in \delta^{-1}(S, b^i)$
	for some $i \ge 0$. 
	Since $6 \notin \delta^{-1}(S, b^i)$
	we also find that $0$ and $3$ are not contained in $\delta^{-1}(S, b^i)$.
	For $T = \delta^{-1}(S, ab^i)$,
	this gives $T = ( \delta^{-1}(S, b^i) \setminus \{ 7  \} ) \cup \{ 6\}$
	and $\delta(T, ab^i) = S$.
	We have $0 \notin \delta^{-1}(T, b^3)$
	since $3 \notin T$ and $3 \in \delta^{-1}(T, b^3)$
	since $6 \in T$.
	We also have $7 \in \delta^{-1}(T, b^3)$ since $1 \in T$
	and $6 \notin \delta^{-1}(T, b^3)$ since $0 \notin T$.
	Thus we find
	\[
	\delta^{-1}(T, ab^3) = ( \delta^{-1}(T, b^3) \setminus \{ 7 \} ) \cup \{ 0,6 \}
	\]
	and $\delta(\delta^{-1}(T, ab^3), ab^3) = T$.
	Since $|\delta^{-1}(T, ab^3)| = |T| + 1 = |S| + 1$, inductively,
	we have $\delta^{-1}(T, ab^3) = \delta(Q, u)$ for some $u \in \Sigma^*$.
	Then we find $\delta(Q, uab^3ab^i) = S$.	
	\end{enumerate}
	
	\begin{figure}[htbp]
		\centering
		\begin{tikzpicture} %[->, >=Stealth, semithick]
			
			\foreach \i in {0,...,8} {
				\node[state] (q\i) at ({40*\i}:3cm) {$\i$};
			}
			
			\foreach \i/\j in {0/1,1/2,2/3,3/4,4/5,5/6,6/7,7/8,8/0} {
				\draw (q\i) edge[dashed] (q\j);
			}
			
			\draw (q0) to (q3);
			\draw (q6) to[bend left=25] (q7);
			\draw (q7) to[bend left=25] (q6);
		\end{tikzpicture}
		\caption[Example completely reachable circular automaton]{Automaton from \cref{syncmax:ex:circular-compl-reach-two}. Solid and dashed edges show the actions 
			of~$a$ and, resp., $b$; self-loop omitted.}
		\label{syncmax:fig:example-cr-circular}
	\end{figure}
\end{example}
